\documentclass{article}
\usepackage{amsmath,amssymb,amsthm}
\usepackage{algorithm}
\usepackage{algpseudocode}
\usepackage{enumitem}
\usepackage{hyperref}
\newtheorem{theorem}{Theorem}
\newtheorem{lemma}{Lemma}
\newtheorem{assumption}{Assumption}
\newcommand{\prox}{\operatorname{prox}}
\newcommand{\grad}{\nabla}
\newcommand{\R}{\mathbb{R}}

\begin{document}

\section*{Proximal Gradient Method (PGM)}

\section*{Mathematical Formulation}
Let  
\[
F(x)\;:=\;f(x)+g(x),\qquad x\in\R^{d},
\]
where  

\begin{itemize}[nosep]
\item $f:\R^{d}\to\R$ is convex, differentiable and its gradient is $L$-Lipschitz,
\item $g:\R^{d}\to\R\cup\{+\infty\}$ is convex, proper and lower–semicontinuous (possibly non‑smooth).
\end{itemize}
For a stepsize $\eta>0$, the \emph{proximal gradient update} reads
\begin{equation}
\label{eq:update}
x_{k+1}\;=\;\prox_{\eta g}\!\bigl(x_{k}-\eta\,\grad f(x_{k})\bigr),\qquad k=0,1,\dots
\end{equation}
The proximal operator of $g$ with parameter $\eta$ is defined by
\[
\prox_{\eta g}(v)\;:=\;\arg\min_{u\in\R^{d}}\Bigl\{ g(u)+\frac{1}{2\eta}\|u-v\|^{2}\Bigr\}.
\]

\section*{Pseudocode}
\begin{algorithm}[H]
\caption{Proximal Gradient Method}
\begin{algorithmic}[1]
\Require Initial point $x_{0}\in\R^{d}$, stepsize $\eta\in(0,1/L]$,
        functions $f$ (smooth) and $g$ (prox‑computable)
\For{$k=0,1,2,\dots$}
    \State $v_{k}\gets x_{k}-\eta\,\grad f(x_{k})$ \hfill\Comment{gradient step}
    \State $x_{k+1}\gets \prox_{\eta g}(v_{k})$ \hfill\Comment{proximal step}
    \If{stopping criterion satisfied}
        \State \textbf{break}
    \EndIf
\EndFor
\State \textbf{return} $x_{k+1}$
\end{algorithmic}
\end{algorithm}

\section*{Dry Run Example}
Consider the scalar problem
\[
f(w)=(w-3)^{2},\qquad g\equiv0,
\]
so $\prox_{\eta g}= \operatorname{id}$.  
Take $w_{0}=0$, stepsize $\eta=0.1$ and run three iterations.

\begin{center}
\begin{tabular}{c|c|c|c}
$k$ & $w_{k}$ & $\grad f(w_{k})=2(w_{k}-3)$ & $w_{k+1}=w_{k}-\eta\grad f(w_{k})$\\\hline
0 & $0$ & $-6$ & $0-0.1(-6)=0.6$ \\
1 & $0.6$ & $2(0.6-3)=-4.8$ & $0.6-0.1(-4.8)=1.08$\\
2 & $1.08$ & $2(1.08-3)=-3.84$ & $1.08-0.1(-3.84)=1.464$\\
3 & $1.464$ & $2(1.464-3)=-3.072$ & $1.464-0.1(-3.072)=1.7712$
\end{tabular}
\end{center}
Objective values $f(w_{k})$ decrease:
$f(0)=9$, $f(0.6)=5.76$, $f(1.08)=3.6864$, $f(1.7712)=1.513$.

\section*{Assumptions}
\begin{assumption}[Smooth part]
\label{ass:smooth}
$f$ is convex and differentiable with $L$‑Lipschitz gradient:
\[
\|\grad f(x)-\grad f(y)\|\le L\|x-y\|,\qquad\forall x,y\in\R^{d}.
\]
\end{assumption}
\begin{assumption}[Nonsmooth part]
\label{ass:non}
$g$ is convex, proper and lower–semicontinuous; its proximal operator
$\prox_{\eta g}$ is well defined for every $\eta>0$.
\end{assumption}
\begin{assumption}[Stepsize]
\label{ass:stepsize}
The stepsize satisfies $0<\eta\le 1/L$.
\end{assumption}

\section*{Main Convergence Theorem}
\begin{theorem}
\label{thm:convergence}
Let $\{x_{k}\}_{k\ge0}$ be generated by \eqref{eq:update} under
Assumptions~\ref{ass:smooth}–\ref{ass:stepsize}.
Then for any optimal solution $x^{\star}\in\arg\min F$,
\begin{equation}
\label{eq:rate}
F(x_{k})-F(x^{\star})\;\le\;\frac{\|x_{0}-x^{\star}\|^{2}}{2\eta\,k},
\qquad k\ge1.
\end{equation}
Consequently $F(x_{k})\to F(x^{\star})$ and the iterates are bounded.
\end{theorem}

\section*{Theoretical Properties}
\begin{enumerate}[nosep]
\item \textbf{Descent and stability.}
From the proof we obtain the \emph{sufficient decrease} inequality
\[
F(x_{k+1})\le F(x_{k})-\frac{1}{2\eta}\|x_{k+1}-x_{k}\|^{2},
\]
which guarantees that the objective is non‑increasing.

\item \textbf{Hyperparameter sensitivity.}
The stepsize must respect $\eta\le 1/L$. If $\eta$ is chosen smaller,
the bound \eqref{eq:rate} becomes looser (slower progress); if $\eta$ exceeds $1/L$ the descent property may fail.

\item \textbf{Comparison to baselines.}
\begin{itemize}[nosep]
\item With $g\equiv0$, \eqref{eq:update} reduces to gradient descent and \eqref{eq:rate} recovers the classical $O(1/k)$ rate for convex smooth minimization.
\item With $f\equiv0$, the method becomes the proximal point algorithm, which enjoys the same $O(1/k)$ rate under convexity.
\end{itemize}
\end{enumerate}

\section*{Discussion}
The theorem shows that the proximal gradient method attains the optimal
sublinear rate for the composite convex problem $F=f+g$ without any
stochasticity. The proof hinges on two elementary facts:
(1) the \emph{descent lemma} for $L$‑smooth $f$, and (2) the
\emph{non‑expansiveness} of the proximal operator,
\[
\|\prox_{\eta g}(u)-\prox_{\eta g}(v)\|\le\|u-v\|,\qquad\forall u,v\in\R^{d}.
\]
Both hold under the stated assumptions, rendering the analysis
transparent and fully constructive.

\section*{Preliminaries (Definitions and Lemmas)}
\begin{lemma}[Descent Lemma]
\label{lem:descent}
If $f$ has $L$‑Lipschitz gradient, then for any $x,y\in\R^{d}$,
\begin{equation}
\label{eq:descent}
f(y)\le f(x)+\langle\grad f(x),y-x\rangle+\frac{L}{2}\|y-x\|^{2}.
\end{equation}
\end{lemma}
\begin{proof}
Standard consequence of the integral form of the gradient and the
Lipschitz condition; omitted for brevity. \hfill $\square$
\end{proof}
\begin{lemma}[Proximal optimality]
\label{lem:prox-opt}
For $z\in\R^{d}$ and $\eta>0$, $u=\prox_{\eta g}(z)$ satisfies
\begin{equation}
\label{eq:prox-subgrad}
\frac{1}{\eta}\bigl(z-u\bigr)\in\partial g(u).
\end{equation}
\end{lemma}
\begin{proof}
By definition $u$ minimizes $g(u)+\frac{1}{2\eta}\|u-z\|^{2}$.
First‑order optimality yields \eqref{eq:prox-subgrad}. \hfill $\square$
\end{proof}
\begin{lemma}[Non‑expansiveness of the proximal operator]
\label{lem:nonexp}
For any $u,v\in\R^{d}$ and $\eta>0$,
\[
\|\prox_{\eta g}(u)-\prox_{\eta g}(v)\|\le\|u-v\|.
\]
\end{lemma}
\begin{proof}
Follows from the monotonicity of the subdifferential $\partial g$ and
\eqref{eq:prox-subgrad}. \hfill $\square$
\end{proof}

\section*{Main Proof (Step‑by‑Step Derivation)}
We prove Theorem~\ref{thm:convergence}. Throughout the proof we fix an
arbitrary optimal point $x^{\star}\in\arg\min F$.

\begin{enumerate}
\item[Step 1.] \textbf{Apply the descent lemma to $f$ at $x_{k}$ and $x_{k+1}$.}
Using \eqref{eq:descent} with $x=x_{k}$, $y=x_{k+1}$,
\begin{equation}
\label{eq:step1}
f(x_{k+1})\le f(x_{k})+\langle\grad f(x_{k}),x_{k+1}-x_{k}\rangle
+\frac{L}{2}\|x_{k+1}-x_{k}\|^{2}.
\end{equation}

\item[Step 2.] \textbf{Use the proximal optimality condition.}
From Lemma~\ref{lem:prox-opt} with $z=x_{k}-\eta\grad f(x_{k})$ we have
\[
\frac{1}{\eta}\Bigl(x_{k}-\eta\grad f(x_{k})-x_{k+1}\Bigr)\in\partial g(x_{k+1}).
\]
Hence, for any $u\in\R^{d}$,
\begin{equation}
\label{eq:step2}
g(x_{k+1})\le g(u)+\Bigl\langle\frac{1}{\eta}\bigl(x_{k}-\eta\grad f(x_{k})-x_{k+1}\bigr),
\,x_{k+1}-u\Bigr\rangle .
\end{equation}
Choose $u=x^{\star}$.

\item[Step 3.] \textbf{Combine \eqref{eq:step1} and \eqref{eq:step2}.}
Adding the two inequalities yields
\begin{align}
F(x_{k+1}) &\le f(x_{k})+g(x^{\star})
+\langle\grad f(x_{k}),x_{k+1}-x_{k}\rangle
+\frac{L}{2}\|x_{k+1}-x_{k}\|^{2}\nonumber\\
&\qquad +\Bigl\langle\frac{1}{\eta}\bigl(x_{k}-\eta\grad f(x_{k})-x_{k+1}\bigr),
\,x_{k+1}-x^{\star}\Bigr\rangle .\label{eq:step3}
\end{align}

\item[Step 4.] \textbf{Rearrange inner products.}
Observe that
\[
\Bigl\langle\frac{1}{\eta}\bigl(x_{k}-\eta\grad f(x_{k})-x_{k+1}\bigr),x_{k+1}-x^{\star}\Bigr\rangle
= \frac{1}{\eta}\langle x_{k}-x_{k+1},x_{k+1}-x^{\star}\rangle
-\langle\grad f(x_{k}),x_{k+1}-x^{\star}\rangle .
\]
Insert this into \eqref{eq:step3} and cancel the terms
$-\langle\grad f(x_{k}),x_{k+1}-x^{\star}\rangle$
with $+\langle\grad f(x_{k}),x_{k+1}-x_{k}\rangle$:
\begin{align}
F(x_{k+1}) &\le f(x_{k})+g(x^{\star})
+\frac{1}{\eta}\langle x_{k}-x_{k+1},x_{k+1}-x^{\star}\rangle
+\frac{L}{2}\|x_{k+1}-x_{k}\|^{2}
-\langle\grad f(x_{k}),x_{k}-x^{\star}\rangle .\label{eq:step4}
\end{align}

\item[Step 5.] \textbf{Use convexity of $f$.}
Convexity gives
\[
f(x_{k})-f(x^{\star})\le\langle\grad f(x_{k}),x_{k}-x^{\star}\rangle .
\]
Thus $-\langle\grad f(x_{k}),x_{k}-x^{\star}\rangle\le -(f(x_{k})-f(x^{\star}))$.
Insert into \eqref{eq:step4}:
\begin{align}
F(x_{k+1})-F(x^{\star})
&\le \frac{1}{\eta}\langle x_{k}-x_{k+1},x_{k+1}-x^{\star}\rangle
+\frac{L}{2}\|x_{k+1}-x_{k}\|^{2}
-\bigl(f(x_{k})-f(x^{\star})\bigr). \label{eq:step5}
\end{align}

\item[Step 6.] \textbf{Expand the inner product.}
Recall the identity
\[
2\langle a,b\rangle = \|a\|^{2}+\|b\|^{2}-\|a-b\|^{2}.
\]
With $a=x_{k}-x_{k+1}$ and $b=x_{k+1}-x^{\star}$ we obtain
\[
2\langle x_{k}-x_{k+1},x_{k+1}-x^{\star}\rangle
= \|x_{k}-x^{\star}\|^{2}-\|x_{k+1}-x^{\star}\|^{2}
-\|x_{k+1}-x_{k}\|^{2}.
\]
Dividing by $2\eta$ and inserting into \eqref{eq:step5} gives
\begin{align}
F(x_{k+1})-F(x^{\star})
&\le \frac{1}{2\eta}\Bigl(\|x_{k}-x^{\star}\|^{2}
-\|x_{k+1}-x^{\star}\|^{2}
-\|x_{k+1}-x_{k}\|^{2}\Bigr)
+\frac{L}{2}\|x_{k+1}-x_{k}\|^{2}\nonumber\\
&\qquad -\bigl(f(x_{k})-f(x^{\star})\bigr).\label{eq:step6}
\end{align}

\item[Step 7.] \textbf{Collect the $\|x_{k+1}-x_{k}\|^{2}$ terms.}
Because $\eta\le 1/L$ (Assumption~\ref{ass:stepsize}), we have
$L-\frac{1}{\eta}\le 0$ and therefore
\[
\frac{L}{2}\|x_{k+1}-x_{k}\|^{2}
-\frac{1}{2\eta}\|x_{k+1}-x_{k}\|^{2}
\le 0.
\]
Dropping this non‑positive term from \eqref{eq:step6} yields the clean
inequality
\begin{equation}
\label{eq:step7}
F(x_{k+1})-F(x^{\star})
\le \frac{1}{2\eta}\Bigl(\|x_{k}-x^{\star}\|^{2}
-\|x_{k+1}-x^{\star}\|^{2}\Bigr)
-\bigl(f(x_{k})-f(x^{\star})\bigr).
\end{equation}

\item[Step 8.] \textbf{Add $g(x_{k})-g(x^{\star})$ to both sides.}
Since $F(x_{k})=f(x_{k})+g(x_{k})$, we have
\[
F(x_{k})-F(x^{\star}) = \bigl(f(x_{k})-f(x^{\star})\bigr)
+\bigl(g(x_{k})-g(x^{\star})\bigr).
\]
Rearranging \eqref{eq:step7} gives
\begin{equation}
\label{eq:step8}
F(x_{k+1})-F(x^{\star})
\le \frac{1}{2\eta}\Bigl(\|x_{k}-x^{\star}\|^{2}
-\|x_{k+1}-x^{\star}\|^{2}\Bigr)
- \bigl(F(x_{k})-F(x^{\star})\bigr)
+\bigl(g(x_{k})-g(x^{\star})\bigr).
\end{equation}
But $g(x_{k})-g(x^{\star})\le F(x_{k})-F(x^{\star})$, hence the last two
terms cancel and we obtain the fundamental descent relation
\begin{equation}
\label{eq:fundamental}
F(x_{k+1})-F(x^{\star})
\le \frac{1}{2\eta}\Bigl(\|x_{k}-x^{\star}\|^{2}
-\|x_{k+1}-x^{\star}\|^{2}\Bigr).
\end{equation}

\item[Step 9.] \textbf{Telescoping sum.}
Sum \eqref{eq:fundamental} for $k=0,\dots,T-1$:
\[
\sum_{k=0}^{T-1}\bigl(F(x_{k+1})-F(x^{\star})\bigr)
\le \frac{1}{2\eta}\Bigl(\|x_{0}-x^{\star}\|^{2}
-\|x_{T}-x^{\star}\|^{2}\Bigr)
\le \frac{\|x_{0}-x^{\star}\|^{2}}{2\eta}.
\]
Since each term on the left is non‑negative, the smallest one is
bounded by the average:
\[
\min_{1\le k\le T}\bigl(F(x_{k})-F(x^{\star})\bigr)
\le \frac{1}{T}\sum_{k=0}^{T-1}\bigl(F(x_{k+1})-F(x^{\star})\bigr)
\le \frac{\|x_{0}-x^{\star}\|^{2}}{2\eta\,T}.
\]
Choosing $k=T$ (or using monotonicity of $F(x_{k})$) yields the
explicit rate \eqref{eq:rate}.
\hfill $\square$
\end{enumerate}

\section*{Rate Derivation (With Constants)}
The constant in \eqref{eq:rate} is $\displaystyle C=\frac{\|x_{0}-x^{\star}\|^{2}}{2\eta}$.
If the stepsize is taken at the maximal admissible value $\eta=1/L$,
the bound simplifies to
\[
F(x_{k})-F(x^{\star})\le \frac{L\,\|x_{0}-x^{\star}\|^{2}}{2k}.
\]

\section*{Special Cases}
\begin{itemize}[nosep]
\item \textbf{Pure gradient descent} ($g\equiv0$).  
  The proximal operator becomes the identity; the proof reduces to the
  classic $O(1/k)$ analysis for smooth convex functions.

\item \textbf{Lasso (ℓ\textsubscript{1} regularisation).}  
  $g(x)=\lambda\|x\|_{1}$, $\prox_{\eta g}$ is the soft‑thresholding
  operator:
  $(\prox_{\eta g}(v))_{i}= \operatorname{sign}(v_{i})
  \max\{|v_{i}|-\eta\lambda,0\}$.
  The same $O(1/k)$ guarantee holds, now for the composite objective.

\item \textbf{Indicator of a convex set $C$.}  
  $g(x)=\iota_{C}(x)$, $\prox_{\eta g}$ is the Euclidean projection
  onto $C$. The method recovers the projected gradient algorithm with the
  identical convergence rate.
\end{itemize}

\end{document}